\documentclass[aspectratio=169,11pt]{beamer}
\usetheme{Madrid}
\usecolortheme{default}
\usepackage[T1]{fontenc}
\setbeamertemplate{navigation symbols}{}
\setbeamertemplate{footline}[frame number]
\usepackage{enumitem}
\usepackage{amsmath,amssymb}
\usepackage{booktabs}
\usepackage{array}

\title{Databricks ML Associate}
\subtitle{High-Yield Next 100 (Set 4)}
\author{Practice Exam Deck}
\date{}

\begin{document}

\begin{frame}
\titlepage
\end{frame}

\begin{frame}{How to Use This Deck}
\begin{itemize}
\item 100 more questions (Q1--Q100), complementing Sets 1 to 3, on topics that commonly appear on the Databricks Machine Learning Associate exam.
\item Each question slide has four options (A--D); the next slide gives the correct answer and why the others are wrong.
\item Topic order: Spark ML transformers; MLflow tracking extras; Delta and streaming; statistics and EDA; ML theory; distributed ML; Hyperopt extras; evaluation; deployment and monitoring; workspace.
\end{itemize}
\end{frame}

\begin{frame}{Exam Domains Covered}
\begin{itemize}
\item \textbf{Databricks Machine Learning:} runtime, AutoML, MLflow, Unity Catalog registry, Feature Engineering client.
\item \textbf{Data Processing:} EDA, imputation, encoding, Spark ML transformers, pandas API on Spark and UDFs.
\item \textbf{Model Development:} Hyperopt, CrossValidator, metrics, bias/variance, leakage.
\item \textbf{Model Deployment:} batch, streaming, real-time serving, monitoring.
\end{itemize}
\end{frame}

% ---------------------------------------------------------------
\begin{frame}{Question 1}
In Spark 3 and later, \texttt{OneHotEncoder} is which type of Spark ML object?
\begin{enumerate}[label=\Alph*.]
\item A model registry client
\item A plain transformer only
\item An estimator that must be fit before \texttt{transform}
\item An evaluator
\end{enumerate}
\end{frame}
\begin{frame}{Answer 1}
\textbf{Correct: C} --- it learns the category sizes during \texttt{fit}.\\
\vspace{0.5em}
\textbf{Wrong:}
\begin{itemize}
\item A: unrelated.
\item B: it needs to be fit in Spark 3+.
\item D: evaluators compute metrics.
\end{itemize}
\end{frame}

% ---------------------------------------------------------------
\begin{frame}{Question 2}
What does \texttt{OneHotEncoder(dropLast=True)} (the default) do for a feature with 5 categories?
\begin{enumerate}[label=\Alph*.]
\item Produces a vector of length 5
\item Produces a vector of length 1
\item Drops the entire column
\item Produces a vector of length 4, dropping the last category
\end{enumerate}
\end{frame}
\begin{frame}{Answer 2}
\textbf{Correct: D} --- the last category is the implicit baseline.\\
\vspace{0.5em}
\textbf{Wrong:}
\begin{itemize}
\item A: that is \texttt{dropLast=False}.
\item B: not the behavior.
\item C: not the behavior.
\end{itemize}
\end{frame}

% ---------------------------------------------------------------
\begin{frame}{Question 3}
By default, which category does \texttt{StringIndexer} assign index 0?
\begin{enumerate}[label=\Alph*.]
\item The alphabetically first label
\item The most frequent label
\item A random label
\item The rarest label
\end{enumerate}
\end{frame}
\begin{frame}{Answer 3}
\textbf{Correct: B} --- the default order is by descending frequency.\\
\vspace{0.5em}
\textbf{Wrong:}
\begin{itemize}
\item A: not the default.
\item C: not the default.
\item D: not the default.
\end{itemize}
\end{frame}

% ---------------------------------------------------------------
\begin{frame}{Question 4}
Which Spark ML class converts predicted numeric indices back to original string labels?
\begin{enumerate}[label=\Alph*.]
\item \texttt{VectorSlicer}
\item \texttt{LabelDecoder}
\item \texttt{IndexToString}
\item \texttt{StringIndexerInverse}
\end{enumerate}
\end{frame}
\begin{frame}{Answer 4}
\textbf{Correct: C} --- it reverses \texttt{StringIndexer}.\\
\vspace{0.5em}
\textbf{Wrong:}
\begin{itemize}
\item A: slices vectors.
\item B: not a Spark ML class.
\item D: not a class.
\end{itemize}
\end{frame}

% ---------------------------------------------------------------
\begin{frame}{Question 5}
Which Spark ML transformer splits a text column into words?
\begin{enumerate}[label=\Alph*.]
\item \texttt{PCA}
\item \texttt{Tokenizer}
\item \texttt{Bucketizer}
\item \texttt{Imputer}
\end{enumerate}
\end{frame}
\begin{frame}{Answer 5}
\textbf{Correct: B} --- it lowercases and splits on whitespace.\\
\vspace{0.5em}
\textbf{Wrong:}
\begin{itemize}
\item A: reduces dimensions.
\item C: bins numbers.
\item D: fills missing values.
\end{itemize}
\end{frame}

% ---------------------------------------------------------------
\begin{frame}{Question 6}
Which pair of Spark ML stages builds TF-IDF features from tokenized text?
\begin{enumerate}[label=\Alph*.]
\item \texttt{Bucketizer} followed by \texttt{Imputer}
\item \texttt{VectorAssembler} followed by \texttt{KMeans}
\item \texttt{HashingTF} (or \texttt{CountVectorizer}) followed by \texttt{IDF}
\item \texttt{StringIndexer} followed by \texttt{PCA}
\end{enumerate}
\end{frame}
\begin{frame}{Answer 6}
\textbf{Correct: C} --- term frequencies are rescaled by inverse document frequency.\\
\vspace{0.5em}
\textbf{Wrong:}
\begin{itemize}
\item A: not TF-IDF.
\item B: not TF-IDF.
\item D: not TF-IDF.
\end{itemize}
\end{frame}

% ---------------------------------------------------------------
\begin{frame}{Question 7}
Which Spark ML transformer scales each row (sample) to unit norm?
\begin{enumerate}[label=\Alph*.]
\item \texttt{StandardScaler}
\item \texttt{Bucketizer}
\item \texttt{MinMaxScaler}
\item \texttt{Normalizer}
\end{enumerate}
\end{frame}
\begin{frame}{Answer 7}
\textbf{Correct: D} --- it normalizes per row, unlike scalers that work per column.\\
\vspace{0.5em}
\textbf{Wrong:}
\begin{itemize}
\item A: scales per feature.
\item B: bins values.
\item C: scales per feature.
\end{itemize}
\end{frame}

% ---------------------------------------------------------------
\begin{frame}{Question 8}
Which scaler preserves sparsity by dividing each feature by its maximum absolute value?
\begin{enumerate}[label=\Alph*.]
\item \texttt{OneHotEncoder}
\item \texttt{MaxAbsScaler}
\item \texttt{Imputer}
\item \texttt{StandardScaler} with centering
\end{enumerate}
\end{frame}
\begin{frame}{Answer 8}
\textbf{Correct: B} --- zeros stay zero.\\
\vspace{0.5em}
\textbf{Wrong:}
\begin{itemize}
\item A: encodes categories.
\item C: fills missing values.
\item D: centering destroys sparsity.
\end{itemize}
\end{frame}

% ---------------------------------------------------------------
\begin{frame}{Question 9}
What does \texttt{RFormula} do in Spark ML?
\begin{enumerate}[label=\Alph*.]
\item Computes ROC AUC
\item Registers a model
\item Builds a features vector and label column from an R-style formula such as \texttt{y \textasciitilde{} a + b}
\item Reads Delta tables
\end{enumerate}
\end{frame}
\begin{frame}{Answer 9}
\textbf{Correct: C} --- it handles indexing and encoding for formula terms.\\
\vspace{0.5em}
\textbf{Wrong:}
\begin{itemize}
\item A: that is an evaluator.
\item B: not its role.
\item D: not its role.
\end{itemize}
\end{frame}

% ---------------------------------------------------------------
\begin{frame}{Question 10}
Which transformer lets you apply a SQL statement to a DataFrame inside a pipeline using \texttt{\_\_THIS\_\_}?
\begin{enumerate}[label=\Alph*.]
\item \texttt{RFormula}
\item \texttt{Bucketizer}
\item \texttt{VectorAssembler}
\item \texttt{SQLTransformer}
\end{enumerate}
\end{frame}
\begin{frame}{Answer 10}
\textbf{Correct: D} --- \texttt{\_\_THIS\_\_} stands for the input DataFrame.\\
\vspace{0.5em}
\textbf{Wrong:}
\begin{itemize}
\item A: uses an R formula.
\item B: bins numbers.
\item C: combines columns.
\end{itemize}
\end{frame}

% ---------------------------------------------------------------
\begin{frame}{Question 11}
Which transformer converts a numeric column to 0.0 or 1.0 using a threshold?
\begin{enumerate}[label=\Alph*.]
\item \texttt{PCA}
\item \texttt{Binarizer}
\item \texttt{Normalizer}
\item \texttt{IndexToString}
\end{enumerate}
\end{frame}
\begin{frame}{Answer 11}
\textbf{Correct: B} --- values above the threshold become 1.0.\\
\vspace{0.5em}
\textbf{Wrong:}
\begin{itemize}
\item A: reduces dimensions.
\item C: scales rows.
\item D: maps indices to labels.
\end{itemize}
\end{frame}

% ---------------------------------------------------------------
\begin{frame}{Question 12}
Which Spark ML component helps select the most relevant features using statistical tests?
\begin{enumerate}[label=\Alph*.]
\item \texttt{KMeans}
\item \texttt{Tokenizer}
\item \texttt{IndexToString}
\item \texttt{UnivariateFeatureSelector} (or \texttt{ChiSqSelector})
\end{enumerate}
\end{frame}
\begin{frame}{Answer 12}
\textbf{Correct: D} --- they score features against the label.\\
\vspace{0.5em}
\textbf{Wrong:}
\begin{itemize}
\item A: clusters rows.
\item B: splits text.
\item C: maps indices.
\end{itemize}
\end{frame}

% ---------------------------------------------------------------
\begin{frame}{Question 13}
What does \texttt{PCA(k=3, inputCol="features", outputCol="pca")} produce after fit and transform?
\begin{enumerate}[label=\Alph*.]
\item A dataset with 3 rows
\item A model with 3 trees
\item A column with the first 3 principal components
\item A column with 3 clusters
\end{enumerate}
\end{frame}
\begin{frame}{Answer 13}
\textbf{Correct: C} --- k sets the number of components.\\
\vspace{0.5em}
\textbf{Wrong:}
\begin{itemize}
\item A: not PCA.
\item B: not PCA.
\item D: that is KMeans.
\end{itemize}
\end{frame}

% ---------------------------------------------------------------
\begin{frame}{Question 14}
Why should \texttt{VectorAssembler} usually be one of the last feature stages in a pipeline?
\begin{enumerate}[label=\Alph*.]
\item It must run before any imputation
\item It registers the pipeline
\item It combines the already processed columns into the single features vector the model needs
\item It trains the model
\end{enumerate}
\end{frame}
\begin{frame}{Answer 14}
\textbf{Correct: C} --- upstream stages create the processed columns first.\\
\vspace{0.5em}
\textbf{Wrong:}
\begin{itemize}
\item A: it needs processed columns.
\item B: not its role.
\item D: it is a transformer.
\end{itemize}
\end{frame}

% ---------------------------------------------------------------
\begin{frame}{Question 15}
Which call saves a matplotlib figure to a run as an image artifact?
\begin{enumerate}[label=\Alph*.]
\item \texttt{mlflow.log\_metric(fig)}
\item \texttt{mlflow.set\_figure(fig)}
\item \texttt{mlflow.log\_param(fig)}
\item \texttt{mlflow.log\_figure(fig, "plot.png")}
\end{enumerate}
\end{frame}
\begin{frame}{Answer 15}
\textbf{Correct: D} --- log\_figure handles matplotlib and plotly figures.\\
\vspace{0.5em}
\textbf{Wrong:}
\begin{itemize}
\item A: metrics are numbers.
\item B: not an API.
\item C: params are small values.
\end{itemize}
\end{frame}

% ---------------------------------------------------------------
\begin{frame}{Question 16}
Which call saves a Python dictionary as a JSON file artifact in a run?
\begin{enumerate}[label=\Alph*.]
\item \texttt{mlflow.log\_dict(d, "config.json")}
\item \texttt{mlflow.log\_params\_file(d)}
\item \texttt{mlflow.set\_dict(d)}
\item \texttt{mlflow.log\_tag(d)}
\end{enumerate}
\end{frame}
\begin{frame}{Answer 16}
\textbf{Correct: A} --- the dict is stored as an artifact file.\\
\vspace{0.5em}
\textbf{Wrong:}
\begin{itemize}
\item B: not an API.
\item C: not an API.
\item D: not an API.
\end{itemize}
\end{frame}

% ---------------------------------------------------------------
\begin{frame}{Question 17}
Which call stores a string of text as an artifact file?
\begin{enumerate}[label=\Alph*.]
\item \texttt{mlflow.log\_text(text, "notes.txt")}
\item \texttt{mlflow.log\_metric(text)}
\item \texttt{mlflow.log\_string(text)}
\item \texttt{mlflow.set\_tag("notes.txt", file)}
\end{enumerate}
\end{frame}
\begin{frame}{Answer 17}
\textbf{Correct: A} --- log\_text writes a text artifact.\\
\vspace{0.5em}
\textbf{Wrong:}
\begin{itemize}
\item B: metrics are numbers.
\item C: not an API.
\item D: tags are not files.
\end{itemize}
\end{frame}

% ---------------------------------------------------------------
\begin{frame}{Question 18}
A filter string needs runs with \texttt{max\_depth} equal to 5. Params are stored as strings. Which \texttt{filter\_string} is correct?
\begin{enumerate}[label=\Alph*.]
\item \texttt{"max\_depth:5"}
\item \texttt{"params.max\_depth = '5'"}
\item \texttt{"metrics.max\_depth = 5"}
\item \texttt{"tags.max\_depth == 5"}
\end{enumerate}
\end{frame}
\begin{frame}{Answer 18}
\textbf{Correct: B} --- param values are compared as quoted strings.\\
\vspace{0.5em}
\textbf{Wrong:}
\begin{itemize}
\item A: not valid syntax.
\item C: max\_depth is a param, not a metric.
\item D: wrong namespace and operator.
\end{itemize}
\end{frame}

% ---------------------------------------------------------------
\begin{frame}{Question 19}
Which filter string finds runs with RMSE below 3?
\begin{enumerate}[label=\Alph*.]
\item \texttt{"tags.rmse < 3"}
\item \texttt{"metrics.rmse < 3"}
\item \texttt{"rmse = 3"}
\item \texttt{"params.rmse < 3"}
\end{enumerate}
\end{frame}
\begin{frame}{Answer 19}
\textbf{Correct: B} --- metrics are compared numerically.\\
\vspace{0.5em}
\textbf{Wrong:}
\begin{itemize}
\item A: tags are strings.
\item C: namespace missing and not a less-than.
\item D: rmse is a metric.
\end{itemize}
\end{frame}

% ---------------------------------------------------------------
\begin{frame}{Question 20}
Which call returns the full history of a metric such as loss over steps for a run?
\begin{enumerate}[label=\Alph*.]
\item \texttt{mlflow.log\_metric\_history()}
\item \texttt{mlflow.get\_metric(run\_id)}
\item \texttt{mlflow.search\_runs(run\_id)}
\item \texttt{MlflowClient().get\_metric\_history(run\_id, "loss")}
\end{enumerate}
\end{frame}
\begin{frame}{Answer 20}
\textbf{Correct: D} --- it returns all logged values with steps.\\
\vspace{0.5em}
\textbf{Wrong:}
\begin{itemize}
\item A: not an API.
\item B: not an API.
\item C: searches runs, not step history.
\end{itemize}
\end{frame}

% ---------------------------------------------------------------
\begin{frame}{Question 21}
Which call downloads a run's artifacts to local storage?
\begin{enumerate}[label=\Alph*.]
\item \texttt{mlflow.artifacts.download\_artifacts(run\_id=..., artifact\_path=...)}
\item \texttt{mlflow.get\_files()}
\item \texttt{mlflow.pull\_artifacts()}
\item \texttt{mlflow.load\_run\_dir()}
\end{enumerate}
\end{frame}
\begin{frame}{Answer 21}
\textbf{Correct: A} --- it copies artifacts locally.\\
\vspace{0.5em}
\textbf{Wrong:}
\begin{itemize}
\item B: not an API.
\item C: not an API.
\item D: not an API.
\end{itemize}
\end{frame}

% ---------------------------------------------------------------
\begin{frame}{Question 22}
Which call lists experiments in the tracking server?
\begin{enumerate}[label=\Alph*.]
\item \texttt{mlflow.search\_runs()}
\item \texttt{mlflow.search\_experiments()}
\item \texttt{mlflow.get\_experiments\_df()}
\item \texttt{mlflow.list\_models()}
\end{enumerate}
\end{frame}
\begin{frame}{Answer 22}
\textbf{Correct: B} --- it returns experiment objects.\\
\vspace{0.5em}
\textbf{Wrong:}
\begin{itemize}
\item A: lists runs.
\item C: not an API.
\item D: not an API.
\end{itemize}
\end{frame}

% ---------------------------------------------------------------
\begin{frame}{Question 23}
Which call gets an experiment object by its path or name?
\begin{enumerate}[label=\Alph*.]
\item \texttt{mlflow.get\_name(experiment)}
\item \texttt{mlflow.find\_experiment(name)}
\item \texttt{mlflow.search\_runs(name)}
\item \texttt{mlflow.get\_experiment\_by\_name(name)}
\end{enumerate}
\end{frame}
\begin{frame}{Answer 23}
\textbf{Correct: D} --- it returns the experiment or None.\\
\vspace{0.5em}
\textbf{Wrong:}
\begin{itemize}
\item A: not an API.
\item B: not an API.
\item C: searches runs.
\end{itemize}
\end{frame}

% ---------------------------------------------------------------
\begin{frame}{Question 24}
What does \texttt{mlflow.log\_input(dataset)} help record?
\begin{enumerate}[label=\Alph*.]
\item The GPU type
\item The dataset used with a run, for lineage
\item The alias
\item The model signature
\end{enumerate}
\end{frame}
\begin{frame}{Answer 24}
\textbf{Correct: B} --- it ties data to the run.\\
\vspace{0.5em}
\textbf{Wrong:}
\begin{itemize}
\item A: separate.
\item C: separate.
\item D: separate.
\end{itemize}
\end{frame}

% ---------------------------------------------------------------
\begin{frame}{Question 25}
What is stored in a run's \texttt{params} as opposed to its \texttt{metrics}?
\begin{enumerate}[label=\Alph*.]
\item They are identical
\item Params are string inputs such as hyperparameters; metrics are numeric results
\item Params hold images; metrics hold text
\item Params are numeric results only
\end{enumerate}
\end{frame}
\begin{frame}{Answer 25}
\textbf{Correct: B} --- they serve different roles.\\
\vspace{0.5em}
\textbf{Wrong:}
\begin{itemize}
\item A: they differ.
\item C: not true.
\item D: that is metrics.
\end{itemize}
\end{frame}

% ---------------------------------------------------------------
\begin{frame}{Question 26}
Which call deletes a run from tracking?
\begin{enumerate}[label=\Alph*.]
\item \texttt{mlflow.clear()}
\item \texttt{mlflow.end\_run(run\_id)}
\item \texttt{mlflow.remove\_metric(run\_id)}
\item \texttt{mlflow.delete\_run(run\_id)}
\end{enumerate}
\end{frame}
\begin{frame}{Answer 26}
\textbf{Correct: D} --- it moves the run to deleted state.\\
\vspace{0.5em}
\textbf{Wrong:}
\begin{itemize}
\item A: not an API.
\item B: ends the active run.
\item C: not an API.
\end{itemize}
\end{frame}

% ---------------------------------------------------------------
\begin{frame}{Question 27}
What does Delta \texttt{CLONE} do?
\begin{enumerate}[label=\Alph*.]
\item Creates a copy of a table (shallow clones reference the source data files)
\item Deletes a table
\item Encrypts the table
\item Trains a model
\end{enumerate}
\end{frame}
\begin{frame}{Answer 27}
\textbf{Correct: A} --- clones are useful for testing and experiments.\\
\vspace{0.5em}
\textbf{Wrong:}
\begin{itemize}
\item B: not what it does.
\item C: not what it does.
\item D: not what it does.
\end{itemize}
\end{frame}

% ---------------------------------------------------------------
\begin{frame}{Question 28}
What does Delta Change Data Feed provide?
\begin{enumerate}[label=\Alph*.]
\item Row-level changes (inserts, updates, deletes) between table versions
\item Faster clusters
\item Automatic feature selection
\item Model serving
\end{enumerate}
\end{frame}
\begin{frame}{Answer 28}
\textbf{Correct: A} --- downstream jobs can consume changes only.\\
\vspace{0.5em}
\textbf{Wrong:}
\begin{itemize}
\item B: not its purpose.
\item C: not its purpose.
\item D: not its purpose.
\end{itemize}
\end{frame}

% ---------------------------------------------------------------
\begin{frame}{Question 29}
Which setting is required for a Structured Streaming write so the query can recover after restart?
\begin{enumerate}[label=\Alph*.]
\item An MLflow run
\item A GPU cluster
\item A model alias
\item A checkpoint location
\end{enumerate}
\end{frame}
\begin{frame}{Answer 29}
\textbf{Correct: D} --- checkpoints store progress.\\
\vspace{0.5em}
\textbf{Wrong:}
\begin{itemize}
\item A: not required.
\item B: not required.
\item C: not required.
\end{itemize}
\end{frame}

% ---------------------------------------------------------------
\begin{frame}{Question 30}
What does \texttt{trigger(availableNow=True)} do in a streaming query?
\begin{enumerate}[label=\Alph*.]
\item Runs forever with no batches
\item Processes all available data and then stops
\item Deletes the checkpoint
\item Re-reads the whole table each second
\end{enumerate}
\end{frame}
\begin{frame}{Answer 30}
\textbf{Correct: B} --- it behaves like an incremental batch.\\
\vspace{0.5em}
\textbf{Wrong:}
\begin{itemize}
\item A: not the behavior.
\item C: not the behavior.
\item D: not the behavior.
\end{itemize}
\end{frame}

% ---------------------------------------------------------------
\begin{frame}{Question 31}
Which call reads a Delta table as a stream?
\begin{enumerate}[label=\Alph*.]
\item \texttt{spark.table("t").stream()}
\item \texttt{spark.read.stream("t")}
\item \texttt{spark.streamTable("t")}
\item \texttt{spark.readStream.table("catalog.schema.t")}
\end{enumerate}
\end{frame}
\begin{frame}{Answer 31}
\textbf{Correct: D} --- readStream creates a streaming DataFrame.\\
\vspace{0.5em}
\textbf{Wrong:}
\begin{itemize}
\item A: not valid.
\item B: not valid.
\item C: not valid.
\end{itemize}
\end{frame}

% ---------------------------------------------------------------
\begin{frame}{Question 32}
Which streaming sink option lets you run batch logic such as scoring on each micro-batch?
\begin{enumerate}[label=\Alph*.]
\item \texttt{foreachBatch}
\item \texttt{mode("overwrite")}
\item \texttt{foreachPartition} on the driver
\item \texttt{coalesce(1)}
\end{enumerate}
\end{frame}
\begin{frame}{Answer 32}
\textbf{Correct: A} --- it exposes each micro-batch as a DataFrame.\\
\vspace{0.5em}
\textbf{Wrong:}
\begin{itemize}
\item B: not a streaming hook.
\item C: not the micro-batch hook.
\item D: not a streaming hook.
\end{itemize}
\end{frame}

% ---------------------------------------------------------------
\begin{frame}{Question 33}
Which output mode writes only new rows added since the last trigger?
\begin{enumerate}[label=\Alph*.]
\item Overwrite
\item Complete
\item Append
\item Update
\end{enumerate}
\end{frame}
\begin{frame}{Answer 33}
\textbf{Correct: C} --- append emits new rows only.\\
\vspace{0.5em}
\textbf{Wrong:}
\begin{itemize}
\item A: not a streaming output mode name.
\item B: rewrites the whole result.
\item D: emits changed rows.
\end{itemize}
\end{frame}

% ---------------------------------------------------------------
\begin{frame}{Question 34}
A feature pipeline must refresh a feature table whenever new rows land in a source Delta table. Which approach fits?
\begin{enumerate}[label=\Alph*.]
\item Manually edit the feature table
\item Recreate the cluster
\item A streaming or scheduled job that computes features and writes them with \texttt{mode="merge"}
\item Overwrite the model
\end{enumerate}
\end{frame}
\begin{frame}{Answer 34}
\textbf{Correct: C} --- merging by key keeps features current.\\
\vspace{0.5em}
\textbf{Wrong:}
\begin{itemize}
\item A: error prone.
\item B: not related.
\item D: not the feature pipeline.
\end{itemize}
\end{frame}

% ---------------------------------------------------------------
\begin{frame}{Question 35}
Which Delta constraint rejects rows that violate a rule such as \texttt{age >= 0}?
\begin{enumerate}[label=\Alph*.]
\item A \texttt{ZORDER} constraint
\item An \texttt{OPTIMIZE} rule
\item A \texttt{CHECK} constraint
\item A \texttt{VACUUM} rule
\end{enumerate}
\end{frame}
\begin{frame}{Answer 35}
\textbf{Correct: C} --- writes that violate it fail.\\
\vspace{0.5em}
\textbf{Wrong:}
\begin{itemize}
\item A: not a constraint.
\item B: not a constraint.
\item D: not a constraint.
\end{itemize}
\end{frame}

% ---------------------------------------------------------------
\begin{frame}{Question 36}
Why does Delta's ACID support matter for ML pipelines?
\begin{enumerate}[label=\Alph*.]
\item It makes models more accurate
\item It disables time travel
\item Concurrent reads and writes see consistent table versions
\item It removes the need for validation
\end{enumerate}
\end{frame}
\begin{frame}{Answer 36}
\textbf{Correct: C} --- training does not read half-written data.\\
\vspace{0.5em}
\textbf{Wrong:}
\begin{itemize}
\item A: not implied.
\item B: it enables time travel.
\item D: not implied.
\end{itemize}
\end{frame}

% ---------------------------------------------------------------
\begin{frame}{Question 37}
A training dataset must be frozen for audit. Which approach is simplest in Delta?
\begin{enumerate}[label=\Alph*.]
\item Record the table version and read it with time travel, or clone the table
\item Delete the table afterward
\item Rely on the latest version
\item Cache it in a notebook
\end{enumerate}
\end{frame}
\begin{frame}{Answer 37}
\textbf{Correct: A} --- both preserve the exact data.\\
\vspace{0.5em}
\textbf{Wrong:}
\begin{itemize}
\item B: loses data.
\item C: changes over time.
\item D: not durable.
\end{itemize}
\end{frame}

% ---------------------------------------------------------------
\begin{frame}{Question 38}
Which write pattern avoids duplicate feature rows when reprocessing the same data?
\begin{enumerate}[label=\Alph*.]
\item Overwrite with an empty DataFrame
\item Append every time
\item Drop the primary key
\item Merge on the primary key
\end{enumerate}
\end{frame}
\begin{frame}{Answer 38}
\textbf{Correct: D} --- upserts replace instead of duplicating.\\
\vspace{0.5em}
\textbf{Wrong:}
\begin{itemize}
\item A: deletes data.
\item B: creates duplicates.
\item C: removes uniqueness.
\end{itemize}
\end{frame}

% ---------------------------------------------------------------
\begin{frame}{Question 39}
A feature is right-skewed. Which statement about mean and median is typical?
\begin{enumerate}[label=\Alph*.]
\item The mean is larger than the median
\item The mean is smaller than the median
\item They are always equal
\item The median is undefined
\end{enumerate}
\end{frame}
\begin{frame}{Answer 39}
\textbf{Correct: A} --- the long right tail pulls the mean up.\\
\vspace{0.5em}
\textbf{Wrong:}
\begin{itemize}
\item B: typical of left skew.
\item C: only for symmetric data.
\item D: it is defined.
\end{itemize}
\end{frame}

% ---------------------------------------------------------------
\begin{frame}{Question 40}
A value is 70, the mean is 50, and the standard deviation is 10. What is its z-score?
\begin{enumerate}[label=\Alph*.]
\item 2.0
\item 0.2
\item 7.0
\item 20
\end{enumerate}
\end{frame}
\begin{frame}{Answer 40}
\textbf{Correct: A} --- (70 - 50) / 10.\\
\vspace{0.5em}
\textbf{Wrong:}
\begin{itemize}
\item B: inverted.
\item C: incorrect.
\item D: forgets to divide.
\end{itemize}
\end{frame}

% ---------------------------------------------------------------
\begin{frame}{Question 41}
Min-max scaling a value of 30 to the range 0 to 1 when the minimum is 10 and the maximum is 50 gives:
\begin{enumerate}[label=\Alph*.]
\item 0.5
\item 0.6
\item 0.3
\item 1.0
\end{enumerate}
\end{frame}
\begin{frame}{Answer 41}
\textbf{Correct: A} --- (30 - 10) / (50 - 10).\\
\vspace{0.5em}
\textbf{Wrong:}
\begin{itemize}
\item B: incorrect.
\item C: incorrect.
\item D: that would be 50.
\end{itemize}
\end{frame}

% ---------------------------------------------------------------
\begin{frame}{Question 42}
What does a box plot show?
\begin{enumerate}[label=\Alph*.]
\item The confusion matrix
\item Model coefficients
\item Median, quartiles, and potential outliers
\item Only the mean
\end{enumerate}
\end{frame}
\begin{frame}{Answer 42}
\textbf{Correct: C} --- the box spans Q1 to Q3.\\
\vspace{0.5em}
\textbf{Wrong:}
\begin{itemize}
\item A: not shown.
\item B: not shown.
\item D: not just the mean.
\end{itemize}
\end{frame}

% ---------------------------------------------------------------
\begin{frame}{Question 43}
Which correlation measure captures monotonic but non-linear relationships?
\begin{enumerate}[label=\Alph*.]
\item R2 only
\item Variance inflation factor
\item Spearman rank correlation
\item Pearson correlation
\end{enumerate}
\end{frame}
\begin{frame}{Answer 43}
\textbf{Correct: C} --- it uses ranks.\\
\vspace{0.5em}
\textbf{Wrong:}
\begin{itemize}
\item A: not a correlation measure.
\item B: detects collinearity.
\item D: measures linear association.
\end{itemize}
\end{frame}

% ---------------------------------------------------------------
\begin{frame}{Question 44}
Correlation between two features is high. What can you conclude?
\begin{enumerate}[label=\Alph*.]
\item They are identical columns
\item They must both be dropped
\item They are associated, but this does not prove causation
\item One causes the other
\end{enumerate}
\end{frame}
\begin{frame}{Answer 44}
\textbf{Correct: C} --- association is not causation.\\
\vspace{0.5em}
\textbf{Wrong:}
\begin{itemize}
\item A: not implied.
\item B: not required.
\item D: not proven.
\end{itemize}
\end{frame}

% ---------------------------------------------------------------
\begin{frame}{Question 45}
A feature has the same value in every row. What should you usually do?
\begin{enumerate}[label=\Alph*.]
\item Drop it because it carries no information
\item One-hot encode it
\item Use it as the key
\item Keep it as the label
\end{enumerate}
\end{frame}
\begin{frame}{Answer 45}
\textbf{Correct: A} --- zero variance features add nothing.\\
\vspace{0.5em}
\textbf{Wrong:}
\begin{itemize}
\item B: adds nothing.
\item C: not appropriate.
\item D: not appropriate.
\end{itemize}
\end{frame}

% ---------------------------------------------------------------
\begin{frame}{Question 46}
Which chart best shows the distribution of a single numeric feature?
\begin{enumerate}[label=\Alph*.]
\item A Sankey chart
\item A pie chart of rows
\item A treemap
\item A histogram
\end{enumerate}
\end{frame}
\begin{frame}{Answer 46}
\textbf{Correct: D} --- it shows frequencies across bins.\\
\vspace{0.5em}
\textbf{Wrong:}
\begin{itemize}
\item A: not a distribution.
\item B: not a distribution.
\item C: not a distribution.
\end{itemize}
\end{frame}

% ---------------------------------------------------------------
\begin{frame}{Question 47}
A scatter plot of a feature against the target is used to:
\begin{enumerate}[label=\Alph*.]
\item Reveal relationships, trends, and outliers
\item Train the model
\item Deploy the endpoint
\item Register the model
\end{enumerate}
\end{frame}
\begin{frame}{Answer 47}
\textbf{Correct: A} --- it is a basic EDA tool.\\
\vspace{0.5em}
\textbf{Wrong:}
\begin{itemize}
\item B: not its purpose.
\item C: not its purpose.
\item D: not its purpose.
\end{itemize}
\end{frame}

% ---------------------------------------------------------------
\begin{frame}{Question 48}
A sample used for EDA must be representative. Which practice helps most?
\begin{enumerate}[label=\Alph*.]
\item Using only rows with no nulls
\item Random or stratified sampling
\item Using only the newest rows
\item Taking the first 1000 rows of a sorted table
\end{enumerate}
\end{frame}
\begin{frame}{Answer 48}
\textbf{Correct: B} --- it avoids bias in the sample.\\
\vspace{0.5em}
\textbf{Wrong:}
\begin{itemize}
\item A: can be biased.
\item C: can be biased.
\item D: can be biased.
\end{itemize}
\end{frame}

% ---------------------------------------------------------------
\begin{frame}{Question 49}
What does a Naive Bayes classifier assume?
\begin{enumerate}[label=\Alph*.]
\item Features are conditionally independent given the class
\item Features are perfectly correlated
\item Labels are continuous
\item Data is time ordered
\end{enumerate}
\end{frame}
\begin{frame}{Answer 49}
\textbf{Correct: A} --- this simplifies probability estimation.\\
\vspace{0.5em}
\textbf{Wrong:}
\begin{itemize}
\item B: opposite.
\item C: classification uses classes.
\item D: not assumed.
\end{itemize}
\end{frame}

% ---------------------------------------------------------------
\begin{frame}{Question 50}
What does a support vector machine try to maximize?
\begin{enumerate}[label=\Alph*.]
\item The margin between classes
\item The learning rate
\item The number of trees
\item The number of clusters
\end{enumerate}
\end{frame}
\begin{frame}{Answer 50}
\textbf{Correct: A} --- it finds the separating boundary with the widest margin.\\
\vspace{0.5em}
\textbf{Wrong:}
\begin{itemize}
\item B: not SVM.
\item C: not SVM.
\item D: not SVM.
\end{itemize}
\end{frame}

% ---------------------------------------------------------------
\begin{frame}{Question 51}
Which criteria are commonly used by decision trees to choose splits for classification?
\begin{enumerate}[label=\Alph*.]
\item Silhouette or elbow
\item VIF or PCA
\item ROC AUC or R2
\item Gini impurity or entropy
\end{enumerate}
\end{frame}
\begin{frame}{Answer 51}
\textbf{Correct: D} --- they measure node impurity.\\
\vspace{0.5em}
\textbf{Wrong:}
\begin{itemize}
\item A: clustering diagnostics.
\item B: not split criteria.
\item C: evaluation metrics.
\end{itemize}
\end{frame}

% ---------------------------------------------------------------
\begin{frame}{Question 52}
What function maps a linear score to a probability in logistic regression?
\begin{enumerate}[label=\Alph*.]
\item ReLU
\item The identity function
\item Min-max scaling
\item The sigmoid (logistic) function
\end{enumerate}
\end{frame}
\begin{frame}{Answer 52}
\textbf{Correct: D} --- it outputs values between 0 and 1.\\
\vspace{0.5em}
\textbf{Wrong:}
\begin{itemize}
\item A: not used for probabilities here.
\item B: would not bound outputs.
\item C: not the link function.
\end{itemize}
\end{frame}

% ---------------------------------------------------------------
\begin{frame}{Question 53}
Which loss is typically optimized by logistic regression?
\begin{enumerate}[label=\Alph*.]
\item MAPE
\item Log loss (cross-entropy)
\item Hinge on distances only
\item Silhouette
\end{enumerate}
\end{frame}
\begin{frame}{Answer 53}
\textbf{Correct: B} --- it penalizes wrong probabilities.\\
\vspace{0.5em}
\textbf{Wrong:}
\begin{itemize}
\item A: regression metric.
\item C: not logistic regression's loss.
\item D: clustering.
\end{itemize}
\end{frame}

% ---------------------------------------------------------------
\begin{frame}{Question 54}
What does k-means require you to specify in advance?
\begin{enumerate}[label=\Alph*.]
\item The tree depth
\item The learning rate only
\item The labels
\item The number of clusters k
\end{enumerate}
\end{frame}
\begin{frame}{Answer 54}
\textbf{Correct: D} --- k is a hyperparameter.\\
\vspace{0.5em}
\textbf{Wrong:}
\begin{itemize}
\item A: not used.
\item B: not required.
\item C: k-means is unsupervised.
\end{itemize}
\end{frame}

% ---------------------------------------------------------------
\begin{frame}{Question 55}
Why is feature scaling important before k-means?
\begin{enumerate}[label=\Alph*.]
\item It uses labels
\item It uses trees
\item It cannot handle numbers
\item It uses distances, so large-range features would dominate
\end{enumerate}
\end{frame}
\begin{frame}{Answer 55}
\textbf{Correct: D} --- scaling balances feature influence.\\
\vspace{0.5em}
\textbf{Wrong:}
\begin{itemize}
\item A: it does not.
\item B: it does not.
\item C: it can.
\end{itemize}
\end{frame}

% ---------------------------------------------------------------
\begin{frame}{Question 56}
In bagging, how is each tree trained?
\begin{enumerate}[label=\Alph*.]
\item On the previous tree's residuals
\item On a single row
\item On the test data
\item On a bootstrap sample of the training data
\end{enumerate}
\end{frame}
\begin{frame}{Answer 56}
\textbf{Correct: D} --- random resampling introduces diversity.\\
\vspace{0.5em}
\textbf{Wrong:}
\begin{itemize}
\item A: that is boosting.
\item B: not bagging.
\item C: invalid.
\end{itemize}
\end{frame}

% ---------------------------------------------------------------
\begin{frame}{Question 57}
What is a typical trade-off between learning rate and number of trees in gradient boosting?
\begin{enumerate}[label=\Alph*.]
\item They are unrelated
\item A smaller rate needs fewer trees
\item A larger rate always generalizes better
\item A smaller learning rate usually needs more trees
\end{enumerate}
\end{frame}
\begin{frame}{Answer 57}
\textbf{Correct: D} --- each tree contributes less.\\
\vspace{0.5em}
\textbf{Wrong:}
\begin{itemize}
\item A: they interact.
\item B: opposite.
\item C: can overfit.
\end{itemize}
\end{frame}

% ---------------------------------------------------------------
\begin{frame}{Question 58}
Which technique creates synthetic minority examples to address imbalance?
\begin{enumerate}[label=\Alph*.]
\item SMOTE
\item Bucketizer
\item VIF
\item PCA
\end{enumerate}
\end{frame}
\begin{frame}{Answer 58}
\textbf{Correct: A} --- apply it to training data only.\\
\vspace{0.5em}
\textbf{Wrong:}
\begin{itemize}
\item B: binning.
\item C: multicollinearity measure.
\item D: dimension reduction.
\end{itemize}
\end{frame}

% ---------------------------------------------------------------
\begin{frame}{Question 59}
Where should oversampling be applied to avoid leakage?
\begin{enumerate}[label=\Alph*.]
\item Only on the training data after the split
\item On the full dataset before splitting
\item On the labels only
\item On the test set only
\end{enumerate}
\end{frame}
\begin{frame}{Answer 59}
\textbf{Correct: A} --- test data must remain untouched.\\
\vspace{0.5em}
\textbf{Wrong:}
\begin{itemize}
\item B: leaks synthetic neighbors.
\item C: invalid.
\item D: invalid.
\end{itemize}
\end{frame}

% ---------------------------------------------------------------
\begin{frame}{Question 60}
What is a dummy baseline model used for?
\begin{enumerate}[label=\Alph*.]
\item Deploying the model
\item Replacing validation
\item Providing a minimum performance level to beat
\item Reducing leakage
\end{enumerate}
\end{frame}
\begin{frame}{Answer 60}
\textbf{Correct: C} --- it shows whether the model adds value.\\
\vspace{0.5em}
\textbf{Wrong:}
\begin{itemize}
\item A: not its purpose.
\item B: not its purpose.
\item D: not its purpose.
\end{itemize}
\end{frame}

% ---------------------------------------------------------------
\begin{frame}{Question 61}
A tree-based importance score can be biased toward which kind of feature?
\begin{enumerate}[label=\Alph*.]
\item High-cardinality or continuous features
\item Constant features
\item Binary features
\item Label columns only
\end{enumerate}
\end{frame}
\begin{frame}{Answer 61}
\textbf{Correct: A} --- they offer many split points.\\
\vspace{0.5em}
\textbf{Wrong:}
\begin{itemize}
\item B: have no splits.
\item C: less likely to be biased upward.
\item D: not relevant.
\end{itemize}
\end{frame}

% ---------------------------------------------------------------
\begin{frame}{Question 62}
What is the main purpose of a train/validation/test split?
\begin{enumerate}[label=\Alph*.]
\item Only the test set is used
\item Train fits the model, validation tunes it, and test gives a final unbiased estimate
\item All three are used to fit weights
\item They are shuffled together at inference
\end{enumerate}
\end{frame}
\begin{frame}{Answer 62}
\textbf{Correct: B} --- each split has a distinct role.\\
\vspace{0.5em}
\textbf{Wrong:}
\begin{itemize}
\item A: need training data.
\item C: validation and test are held out.
\item D: not a purpose.
\end{itemize}
\end{frame}

% ---------------------------------------------------------------
\begin{frame}{Question 63}
In linear regression, what do residuals represent?
\begin{enumerate}[label=\Alph*.]
\item The regularization terms
\item The learning rates
\item The differences between actual and predicted values
\item The feature means
\end{enumerate}
\end{frame}
\begin{frame}{Answer 63}
\textbf{Correct: C} --- patterns in residuals indicate model problems.\\
\vspace{0.5em}
\textbf{Wrong:}
\begin{itemize}
\item A: not residuals.
\item B: not residuals.
\item D: not residuals.
\end{itemize}
\end{frame}

% ---------------------------------------------------------------
\begin{frame}{Question 64}
Which statement about stacking is true?
\begin{enumerate}[label=\Alph*.]
\item It only works for clustering
\item A meta-model learns to combine predictions of base models
\item It trains one tree on all labels
\item It removes the need for validation
\end{enumerate}
\end{frame}
\begin{frame}{Answer 64}
\textbf{Correct: B} --- it is an ensemble technique.\\
\vspace{0.5em}
\textbf{Wrong:}
\begin{itemize}
\item A: not true.
\item C: not stacking.
\item D: not true.
\end{itemize}
\end{frame}

% ---------------------------------------------------------------
\begin{frame}{Question 65}
A dataset does not fit in the memory of one machine and you need to train a model. Which choice scales out?
\begin{enumerate}[label=\Alph*.]
\item \texttt{df.toPandas()} then fit
\item scikit-learn on the driver
\item Increasing \texttt{maxIter} only
\item Spark ML or a distributed library such as distributed XGBoost
\end{enumerate}
\end{frame}
\begin{frame}{Answer 65}
\textbf{Correct: D} --- they train across workers.\\
\vspace{0.5em}
\textbf{Wrong:}
\begin{itemize}
\item A: collects to the driver.
\item B: limited to one machine.
\item C: does not fix memory.
\end{itemize}
\end{frame}

% ---------------------------------------------------------------
\begin{frame}{Question 66}
What is data parallelism in distributed training?
\begin{enumerate}[label=\Alph*.]
\item Workers train different unrelated models
\item Each worker trains on a shard of data and gradients or updates are synchronized
\item Only the driver trains
\item Each worker holds a different dataset type
\end{enumerate}
\end{frame}
\begin{frame}{Answer 66}
\textbf{Correct: B} --- the model is replicated across workers.\\
\vspace{0.5em}
\textbf{Wrong:}
\begin{itemize}
\item A: that is a different pattern.
\item C: not parallel.
\item D: not the definition.
\end{itemize}
\end{frame}

% ---------------------------------------------------------------
\begin{frame}{Question 67}
Which Databricks-supported tools distribute deep learning training across GPUs or nodes?
\begin{enumerate}[label=\Alph*.]
\item \texttt{VACUUM}
\item \texttt{TorchDistributor} or Horovod
\item \texttt{OneHotEncoder}
\item \texttt{dbutils.fs}
\end{enumerate}
\end{frame}
\begin{frame}{Answer 67}
\textbf{Correct: B} --- they coordinate distributed training.\\
\vspace{0.5em}
\textbf{Wrong:}
\begin{itemize}
\item A: Delta maintenance.
\item C: feature encoder.
\item D: file utility.
\end{itemize}
\end{frame}

% ---------------------------------------------------------------
\begin{frame}{Question 68}
A team wants to run many small independent model fits (one per product) quickly. Which is a good approach?
\begin{enumerate}[label=\Alph*.]
\item Group by product and train each model in parallel with \texttt{applyInPandas}
\item One huge model on the driver for all products
\item Use \texttt{coalesce(1)}
\item Collect each product to the driver sequentially
\end{enumerate}
\end{frame}
\begin{frame}{Answer 68}
\textbf{Correct: A} --- each group trains on a worker.\\
\vspace{0.5em}
\textbf{Wrong:}
\begin{itemize}
\item B: may not fit the need.
\item C: removes parallelism.
\item D: slow.
\end{itemize}
\end{frame}

% ---------------------------------------------------------------
\begin{frame}{Question 69}
Which action is a safe way to explore a huge DataFrame in pandas?
\begin{enumerate}[label=\Alph*.]
\item Call \texttt{collect()} on everything
\item Take a random sample with \texttt{df.sample(fraction)} and then \texttt{toPandas()}
\item Call \texttt{repartition(1)}
\item Call \texttt{toPandas()} on the full DataFrame
\end{enumerate}
\end{frame}
\begin{frame}{Answer 69}
\textbf{Correct: B} --- sampling avoids driver overflow.\\
\vspace{0.5em}
\textbf{Wrong:}
\begin{itemize}
\item A: may crash the driver.
\item C: does not reduce data.
\item D: may crash the driver.
\end{itemize}
\end{frame}

% ---------------------------------------------------------------
\begin{frame}{Question 70}
Which scaling approach helps when Spark joins create too many small partitions or very large shuffles?
\begin{enumerate}[label=\Alph*.]
\item Convert everything to CSV
\item Tune partitioning and use broadcast joins for small tables
\item Use more registered models
\item Disable caching
\end{enumerate}
\end{frame}
\begin{frame}{Answer 70}
\textbf{Correct: B} --- this reduces shuffle cost.\\
\vspace{0.5em}
\textbf{Wrong:}
\begin{itemize}
\item A: not helpful.
\item C: not related.
\item D: not helpful.
\end{itemize}
\end{frame}

% ---------------------------------------------------------------
\begin{frame}{Question 71}
What does a Single Node cluster lack compared to a multi-node cluster?
\begin{enumerate}[label=\Alph*.]
\item MLflow support
\item Notebook support
\item Worker nodes for distributing Spark tasks
\item A Python runtime
\end{enumerate}
\end{frame}
\begin{frame}{Answer 71}
\textbf{Correct: C} --- only the driver is used.\\
\vspace{0.5em}
\textbf{Wrong:}
\begin{itemize}
\item A: it supports MLflow.
\item B: it supports notebooks.
\item D: it has Python.
\end{itemize}
\end{frame}

% ---------------------------------------------------------------
\begin{frame}{Question 72}
Which resource is most important for running \texttt{toPandas()} on a moderately large result?
\begin{enumerate}[label=\Alph*.]
\item Model alias count
\item Git folder size
\item Driver memory
\item Number of workers
\end{enumerate}
\end{frame}
\begin{frame}{Answer 72}
\textbf{Correct: C} --- the data is collected to the driver.\\
\vspace{0.5em}
\textbf{Wrong:}
\begin{itemize}
\item A: unrelated.
\item B: unrelated.
\item D: does not hold collected data.
\end{itemize}
\end{frame}

% ---------------------------------------------------------------
\begin{frame}{Question 73}
Which \texttt{fmin} argument stops the search after a number of seconds?
\begin{enumerate}[label=\Alph*.]
\item \texttt{max\_depth}
\item \texttt{timeout}
\item \texttt{parallelism}
\item \texttt{numFolds}
\end{enumerate}
\end{frame}
\begin{frame}{Answer 73}
\textbf{Correct: B} --- it bounds total search time.\\
\vspace{0.5em}
\textbf{Wrong:}
\begin{itemize}
\item A: a model hyperparameter.
\item C: SparkTrials argument.
\item D: Spark ML CV.
\end{itemize}
\end{frame}

% ---------------------------------------------------------------
\begin{frame}{Question 74}
Which feature lets Hyperopt stop when the loss has not improved for several trials?
\begin{enumerate}[label=\Alph*.]
\item \texttt{STATUS\_FAIL}
\item \texttt{space\_eval}
\item \texttt{hp.choice} only
\item An early stopping function such as \texttt{no\_progress\_loss}
\end{enumerate}
\end{frame}
\begin{frame}{Answer 74}
\textbf{Correct: D} --- passed as \texttt{early\_stop\_fn}.\\
\vspace{0.5em}
\textbf{Wrong:}
\begin{itemize}
\item A: marks a failed trial.
\item B: decodes parameters.
\item C: defines a space.
\end{itemize}
\end{frame}

% ---------------------------------------------------------------
\begin{frame}{Question 75}
How can you make a Hyperopt search reproducible?
\begin{enumerate}[label=\Alph*.]
\item Disable MLflow
\item Increase \texttt{max\_evals}
\item Pass a seeded random state through \texttt{rstate} (for example \texttt{np.random.default\_rng(42)})
\item Use more workers
\end{enumerate}
\end{frame}
\begin{frame}{Answer 75}
\textbf{Correct: C} --- the seed fixes sampling.\\
\vspace{0.5em}
\textbf{Wrong:}
\begin{itemize}
\item A: does not fix randomness.
\item B: does not fix randomness.
\item D: does not fix randomness.
\end{itemize}
\end{frame}

% ---------------------------------------------------------------
\begin{frame}{Question 76}
How can a single Hyperopt search compare different model types such as RandomForest and GradientBoosting?
\begin{enumerate}[label=\Alph*.]
\item Use \texttt{hp.choice} with nested search spaces for each model type
\item Run \texttt{fmin} on the test set
\item Use \texttt{hp.uniform} only
\item Use \texttt{df.randomSplit}
\end{enumerate}
\end{frame}
\begin{frame}{Answer 76}
\textbf{Correct: A} --- the choice selects which sub-space is sampled.\\
\vspace{0.5em}
\textbf{Wrong:}
\begin{itemize}
\item B: invalid.
\item C: cannot pick model types.
\item D: unrelated.
\end{itemize}
\end{frame}

% ---------------------------------------------------------------
\begin{frame}{Question 77}
What does returning \texttt{\{"loss": None, "status": STATUS\_FAIL\}} from an objective indicate?
\begin{enumerate}[label=\Alph*.]
\item That the trial was best
\item That the model is registered
\item That the trial failed
\item That tuning is complete
\end{enumerate}
\end{frame}
\begin{frame}{Answer 77}
\textbf{Correct: C} --- the status marks the trial outcome.\\
\vspace{0.5em}
\textbf{Wrong:}
\begin{itemize}
\item A: no.
\item B: no.
\item D: no.
\end{itemize}
\end{frame}

% ---------------------------------------------------------------
\begin{frame}{Question 78}
Which method of Hyperopt \texttt{Trials} gives the dictionary describing the best trial?
\begin{enumerate}[label=\Alph*.]
\item \texttt{trials.winner}
\item \texttt{trials.top()}
\item \texttt{trials.predict()}
\item \texttt{trials.best\_trial}
\end{enumerate}
\end{frame}
\begin{frame}{Answer 78}
\textbf{Correct: D} --- it holds the result and misc info for the best trial.\\
\vspace{0.5em}
\textbf{Wrong:}
\begin{itemize}
\item A: not an attribute.
\item B: not an attribute.
\item C: not an attribute.
\end{itemize}
\end{frame}

% ---------------------------------------------------------------
\begin{frame}{Question 79}
What are the axes of an ROC curve?
\begin{enumerate}[label=\Alph*.]
\item Loss versus epoch
\item True positive rate versus false positive rate
\item RMSE versus R2
\item Precision versus recall
\end{enumerate}
\end{frame}
\begin{frame}{Answer 79}
\textbf{Correct: B} --- points are computed across thresholds.\\
\vspace{0.5em}
\textbf{Wrong:}
\begin{itemize}
\item A: a learning curve.
\item C: not an ROC curve.
\item D: that is the PR curve.
\end{itemize}
\end{frame}

% ---------------------------------------------------------------
\begin{frame}{Question 80}
What does an AUC of 0.5 indicate?
\begin{enumerate}[label=\Alph*.]
\item A perfect classifier
\item Performance no better than random guessing
\item A perfectly wrong classifier
\item A calibrated model
\end{enumerate}
\end{frame}
\begin{frame}{Answer 80}
\textbf{Correct: B} --- the ROC curve follows the diagonal.\\
\vspace{0.5em}
\textbf{Wrong:}
\begin{itemize}
\item A: that is 1.0.
\item C: that would be 0.0.
\item D: not implied.
\end{itemize}
\end{frame}

% ---------------------------------------------------------------
\begin{frame}{Question 81}
What is specificity?
\begin{enumerate}[label=\Alph*.]
\item TN / (TN + FP)
\item TP / (TP + FN)
\item TP / (TP + FP)
\item FP / (FP + TN)
\end{enumerate}
\end{frame}
\begin{frame}{Answer 81}
\textbf{Correct: A} --- the true negative rate.\\
\vspace{0.5em}
\textbf{Wrong:}
\begin{itemize}
\item B: that is recall.
\item C: that is precision.
\item D: that is the false positive rate.
\end{itemize}
\end{frame}

% ---------------------------------------------------------------
\begin{frame}{Question 82}
What does balanced accuracy compute?
\begin{enumerate}[label=\Alph*.]
\item The sum of all errors
\item The AUC
\item The maximum class accuracy
\item The average of recall across classes
\end{enumerate}
\end{frame}
\begin{frame}{Answer 82}
\textbf{Correct: D} --- it corrects for class imbalance.\\
\vspace{0.5em}
\textbf{Wrong:}
\begin{itemize}
\item A: not its definition.
\item B: different metric.
\item C: not its definition.
\end{itemize}
\end{frame}

% ---------------------------------------------------------------
\begin{frame}{Question 83}
When is RMSLE preferred over RMSE?
\begin{enumerate}[label=\Alph*.]
\item When targets are negative
\item When relative errors matter and targets span orders of magnitude
\item When classifying images
\item When labels are binary
\end{enumerate}
\end{frame}
\begin{frame}{Answer 83}
\textbf{Correct: B} --- it uses logs of target values.\\
\vspace{0.5em}
\textbf{Wrong:}
\begin{itemize}
\item A: log is not defined.
\item C: not a classification metric.
\item D: not a classification metric.
\end{itemize}
\end{frame}

% ---------------------------------------------------------------
\begin{frame}{Question 84}
For a dataset with 2 percent positives, what is the baseline of the precision-recall curve (random classifier precision)?
\begin{enumerate}[label=\Alph*.]
\item 0.98
\item 0.5
\item About 0.02
\item 1.0
\end{enumerate}
\end{frame}
\begin{frame}{Answer 84}
\textbf{Correct: C} --- random precision equals the positive prevalence.\\
\vspace{0.5em}
\textbf{Wrong:}
\begin{itemize}
\item A: that is the negative rate.
\item B: that is for ROC AUC.
\item D: that is perfect.
\end{itemize}
\end{frame}

% ---------------------------------------------------------------
\begin{frame}{Question 85}
A missed fraud case costs much more than a false alarm. How should the decision threshold be set?
\begin{enumerate}[label=\Alph*.]
\item Remove the threshold
\item Raise the threshold to maximize precision only
\item Lower the threshold to favor recall, guided by the business costs
\item Use 0.5 always
\end{enumerate}
\end{frame}
\begin{frame}{Answer 85}
\textbf{Correct: C} --- cost-sensitive thresholding aligns the model with costs.\\
\vspace{0.5em}
\textbf{Wrong:}
\begin{itemize}
\item A: needed for labels.
\item B: misses more fraud.
\item D: ignores costs.
\end{itemize}
\end{frame}

% ---------------------------------------------------------------
\begin{frame}{Question 86}
10-fold cross-validation trains each model on what share of the data?
\begin{enumerate}[label=\Alph*.]
\item 50 percent
\item 90 percent
\item 10 percent
\item 100 percent
\end{enumerate}
\end{frame}
\begin{frame}{Answer 86}
\textbf{Correct: B} --- 9 of 10 folds are used for training.\\
\vspace{0.5em}
\textbf{Wrong:}
\begin{itemize}
\item A: incorrect.
\item C: that is the validation fold.
\item D: would include validation data.
\end{itemize}
\end{frame}

% ---------------------------------------------------------------
\begin{frame}{Question 87}
Which Lakehouse Monitoring profile type fits a table of model inputs, predictions, and labels?
\begin{enumerate}[label=\Alph*.]
\item Delta profile
\item Time series profile only
\item Inference profile
\item Snapshot profile
\end{enumerate}
\end{frame}
\begin{frame}{Answer 87}
\textbf{Correct: C} --- it is designed for ML inference logs.\\
\vspace{0.5em}
\textbf{Wrong:}
\begin{itemize}
\item A: not a profile type.
\item B: does not include model quality.
\item D: is for non-time-series tables.
\end{itemize}
\end{frame}

% ---------------------------------------------------------------
\begin{frame}{Question 88}
What does Lakehouse Monitoring create when a monitor is set up?
\begin{enumerate}[label=\Alph*.]
\item A new Git repo
\item A new cluster policy
\item A new model alias
\item Metric tables (profile and drift) and an auto-generated dashboard
\end{enumerate}
\end{frame}
\begin{frame}{Answer 88}
\textbf{Correct: D} --- they summarize statistics and drift.\\
\vspace{0.5em}
\textbf{Wrong:}
\begin{itemize}
\item A: not created.
\item B: not created.
\item C: not created.
\end{itemize}
\end{frame}

% ---------------------------------------------------------------
\begin{frame}{Question 89}
What is a baseline table used for in drift monitoring?
\begin{enumerate}[label=\Alph*.]
\item Comparing current data against reference data such as training data
\item Defining aliases
\item Logging job failures
\item Storing model weights
\end{enumerate}
\end{frame}
\begin{frame}{Answer 89}
\textbf{Correct: A} --- drift is measured relative to the baseline.\\
\vspace{0.5em}
\textbf{Wrong:}
\begin{itemize}
\item B: not its purpose.
\item C: not its purpose.
\item D: not its purpose.
\end{itemize}
\end{frame}

% ---------------------------------------------------------------
\begin{frame}{Question 90}
A monitor flags a large drift in a key feature. What is a sensible next step?
\begin{enumerate}[label=\Alph*.]
\item Delete the monitor
\item Increase the endpoint size
\item Investigate the cause and consider retraining or fixing the upstream data
\item Ignore it
\end{enumerate}
\end{frame}
\begin{frame}{Answer 90}
\textbf{Correct: C} --- drift may be real change or a pipeline bug.\\
\vspace{0.5em}
\textbf{Wrong:}
\begin{itemize}
\item A: removes visibility.
\item B: does not fix drift.
\item D: may degrade quality.
\end{itemize}
\end{frame}

% ---------------------------------------------------------------
\begin{frame}{Question 91}
What is a canary deployment?
\begin{enumerate}[label=\Alph*.]
\item Deleting the old model
\item Training on production traffic only
\item Deploying only on weekends
\item Sending a small share of traffic to a new version before full rollout
\end{enumerate}
\end{frame}
\begin{frame}{Answer 91}
\textbf{Correct: D} --- it limits the blast radius of problems.\\
\vspace{0.5em}
\textbf{Wrong:}
\begin{itemize}
\item A: not canary.
\item B: not canary.
\item C: not canary.
\end{itemize}
\end{frame}

% ---------------------------------------------------------------
\begin{frame}{Question 92}
What is shadow deployment?
\begin{enumerate}[label=\Alph*.]
\item Disabling the old model
\item Running a new model on live traffic without using its predictions for decisions
\item Training offline only
\item Serving the new model to all users
\end{enumerate}
\end{frame}
\begin{frame}{Answer 92}
\textbf{Correct: B} --- its outputs are logged and compared.\\
\vspace{0.5em}
\textbf{Wrong:}
\begin{itemize}
\item A: not shadow.
\item C: not shadow.
\item D: not shadow.
\end{itemize}
\end{frame}

% ---------------------------------------------------------------
\begin{frame}{Question 93}
A custom pyfunc model needs to load a vectorizer file once at startup. Which method is designed for this?
\begin{enumerate}[label=\Alph*.]
\item \texttt{\_\_del\_\_}
\item \texttt{start\_run()}
\item \texttt{load\_context(self, context)}
\item \texttt{predict()} on every call
\end{enumerate}
\end{frame}
\begin{frame}{Answer 93}
\textbf{Correct: C} --- it runs when the model is loaded.\\
\vspace{0.5em}
\textbf{Wrong:}
\begin{itemize}
\item A: not designed for loading.
\item B: not a pyfunc method.
\item D: reloading is wasteful.
\end{itemize}
\end{frame}

% ---------------------------------------------------------------
\begin{frame}{Question 94}
Which artifact do you pass to a pyfunc model so it can access extra files at load time?
\begin{enumerate}[label=\Alph*.]
\item \texttt{tags=\{...\}}
\item \texttt{artifacts=\{...\}} in \texttt{mlflow.pyfunc.log\_model}
\item \texttt{conda\_env=\{...\}}
\item \texttt{signature=\{...\}}
\end{enumerate}
\end{frame}
\begin{frame}{Answer 94}
\textbf{Correct: B} --- they are available through \texttt{context.artifacts}.\\
\vspace{0.5em}
\textbf{Wrong:}
\begin{itemize}
\item A: labels the model.
\item C: defines the environment only.
\item D: defines schema.
\end{itemize}
\end{frame}

% ---------------------------------------------------------------
\begin{frame}{Question 95}
A model endpoint is returning errors after a library update in the training notebook. What is the most likely issue?
\begin{enumerate}[label=\Alph*.]
\item The logged model's dependencies do not match the serving environment or code expectations
\item The feature table is a Delta table
\item The endpoint has too many aliases
\item The experiment name changed
\end{enumerate}
\end{frame}
\begin{frame}{Answer 95}
\textbf{Correct: A} --- dependencies should be captured with the model.\\
\vspace{0.5em}
\textbf{Wrong:}
\begin{itemize}
\item B: not the cause.
\item C: not related.
\item D: not related.
\end{itemize}
\end{frame}

% ---------------------------------------------------------------
\begin{frame}{Question 96}
Which option reduces cluster start-up time for frequent jobs?
\begin{enumerate}[label=\Alph*.]
\item Turning off MLflow
\item Running \texttt{VACUUM} more often
\item Instance pools or serverless compute
\item Increasing \texttt{maxDepth}
\end{enumerate}
\end{frame}
\begin{frame}{Answer 96}
\textbf{Correct: C} --- pools keep warm instances ready.\\
\vspace{0.5em}
\textbf{Wrong:}
\begin{itemize}
\item A: not related.
\item B: not related.
\item D: a model hyperparameter.
\end{itemize}
\end{frame}

% ---------------------------------------------------------------
\begin{frame}{Question 97}
Which command exits a notebook and returns a value to the caller (for example in a job or \texttt{notebook.run})?
\begin{enumerate}[label=\Alph*.]
\item \texttt{dbutils.notebook.exit(value)}
\item \texttt{dbutils.fs.rm(value)}
\item \texttt{mlflow.end\_run(value)}
\item \texttt{spark.stop(value)}
\end{enumerate}
\end{frame}
\begin{frame}{Answer 97}
\textbf{Correct: A} --- it ends the notebook with a result string.\\
\vspace{0.5em}
\textbf{Wrong:}
\begin{itemize}
\item B: deletes files.
\item C: ends a run.
\item D: stops Spark.
\end{itemize}
\end{frame}

% ---------------------------------------------------------------
\begin{frame}{Question 98}
Which cell magic runs SQL in a Python notebook?
\begin{enumerate}[label=\Alph*.]
\item \texttt{\%scala} only
\item \texttt{\%sql}
\item \texttt{\%md}
\item \texttt{\%sh}
\end{enumerate}
\end{frame}
\begin{frame}{Answer 98}
\textbf{Correct: B} --- it switches the cell language to SQL.\\
\vspace{0.5em}
\textbf{Wrong:}
\begin{itemize}
\item A: Scala.
\item C: markdown.
\item D: shell.
\end{itemize}
\end{frame}

% ---------------------------------------------------------------
\begin{frame}{Question 99}
Which permission setting controls who can view or edit an MLflow experiment?
\begin{enumerate}[label=\Alph*.]
\item Delta constraints
\item Notebook themes
\item Cluster policies
\item Experiment permissions (ACLs)
\end{enumerate}
\end{frame}
\begin{frame}{Answer 99}
\textbf{Correct: D} --- access is managed per experiment.\\
\vspace{0.5em}
\textbf{Wrong:}
\begin{itemize}
\item A: limit data values.
\item B: cosmetic.
\item C: limit configs.
\end{itemize}
\end{frame}

% ---------------------------------------------------------------
\begin{frame}{Question 100}
Which feature lets teams share governed, reusable SQL functions and models under the same catalog namespace as tables?
\begin{enumerate}[label=\Alph*.]
\item DBFS root
\item Cluster init scripts
\item Notebook revision history
\item Unity Catalog schemas that hold tables, functions, volumes, and models
\end{enumerate}
\end{frame}
\begin{frame}{Answer 100}
\textbf{Correct: D} --- objects live in the same three-level namespace.\\
\vspace{0.5em}
\textbf{Wrong:}
\begin{itemize}
\item A: ungoverned storage.
\item B: not data governance.
\item C: not governance.
\end{itemize}
\end{frame}

\end{document}